%% ============================================================
\section{Generalizability: MLE-Bench and Parameter Golf}
\label{sec:generalizability}


To test whether the discovery loop transfers beyond ADRS, we evaluate \sys{} unmodified across a diverse set of six tasks selected for their rigorous demands and relevance to current AI research. 
First, we select five MLE-Bench~\citep{chan2024mle} Kaggle competitions spanning medical imaging, fine-grained recognition, and 3D perception. 
We target Medium and High difficulty tiers to ensure sufficient task complexity. 
Second, to assess adaptability in a novel, live environment, we evaluate \sys{} on the Parameter Golf competition~\citep{openai2026parametergolf}.
This competition requires training the highest-performing language model under strict size and performance constraints.
Both systems are provided with a knowledge base of official leaderboard solutions up to a cutoff date of April~27, 2026, when the SOTA score was 1.0611. The leaderboard has since advanced beyond this mark.
Full task details and evaluation setups are provided in Appendix~\ref{app:mle-pg-eval}.


\begin{table}[h]
\centering
\small
\caption{Comparison of solver performance across five MLE-Bench tasks and Parameter Golf. ``Above/Below Median'' denotes outperforming or underperforming the median participant. MLE-Bench medals (Gold, Silver, Bronze) represent simulated private leaderboard standings. ``SOTA'' indicates top-1 performance on the Parameter Golf Leaderboard (as of 2026-04-27). }
\label{tab:generalizability}
\begin{tabular}{l c c >{\centering\arraybackslash}m{3cm} c >{\centering\arraybackslash}m{3cm}}
    \toprule
    \multirow{2}{*}{\textbf{Task}} & \multirow{2}{*}{\textbf{Dir.}} & \multicolumn{2}{c}{\textbf{DeepScientist}} & \multicolumn{2}{c}{\textbf{\sys{}}} \\ \cmidrule(lr){3-4} \cmidrule(lr){5-6} 
    & & \textbf{Score} & \textbf{Highlight} & \textbf{Score} & \textbf{Highlight} \\ 
    \midrule
    3D Object Detection & $\uparrow$ & 0.0000 & Below Median & 0.1763 & Gold Medal \\
    AI4Code & $\uparrow$ & 0.6964 & Below Median & 0.8356 & Above Median \\
    iMet 2020 FGVC7 & $\uparrow$ & 0.6804 & Silver Medal & 0.6791 & Silver Medal \\
    RSNA Brain Tumor & $\uparrow$ & 0.6377 & Gold Medal & 0.6518 & Gold Medal \\ 
    iNaturalist 2019 FGVC6 & $\downarrow$ & 0.2158 & Silver Medal & 0.2445 & Silver Medal \\ \midrule
    Parameter Golf & $\downarrow$ & Invalid & Size limit exceeded & 1.0600 & SOTA (Constraints met) \\
    \bottomrule
\end{tabular}
\end{table}

\paragraph{Solver performance generalizes and demonstrates robustness.}
As shown in Table \ref{tab:generalizability}, \sys{} demonstrates strong generalization across diverse domains, difficulty levels, and strict constraints. 
%Crucially, \sys{} succeeds in environments where the baseline completely fails. 
On the High difficulty MLE-Bench tasks, \sys{} earns two Gold Medals (RSNA Brain Tumor and 3D Object Detection), notably solving the 3D Object Detection task with a Gold Medal score while DeepScientist failed entirely (scoring 0.0000). 
On the Medium difficulty tasks, \sys{} secures Silver Medals on both iMet 2020 and iNaturalist 2019—remaining highly competitive with DeepScientist—and achieves an Above Median standing on AI4Code, marking a significant improvement. 
On the Parameter Golf LLM training task---an entirely different domain from ADRS---\sys{} further demonstrates adaptability. 
While the baseline fails to produce a valid submission due to exceeding the 16MB artifact size limit, \sys{} adheres to all constraints and achieves state-of-the-art performance with a score of 1.0600.


\begin{wrapfigure}{r}{0.5\textwidth}
  \centering
  \includegraphics[width=\linewidth]{figures/novel-idea-parameter-golf.png}
  \caption{Overview of the novel ideas generated by \sys{} for Parameter Golf. Key algorithmic innovations include the Hessian-diagonal-weighted SVD initialization and the GPTQ-driven alternating-least-squares (ALS) refinement loop.}
  \label{fig:scientistone_pipeline_parameter_golf}
\end{wrapfigure}

\paragraph{Solution novelty comparison on Parameter Golf.}
Both \sys{} and DeepScientist were provided with the same prior-art reference and achieved superficially similar numerical improvements. 
However, the systems arrived at these outcomes through fundamentally different approaches. 
\sys{} demonstrated genuine research capability by introducing novel algorithmic techniques to the quantization block: specifically, a Hessian-diagonal-weighted SVD initialization and an alternating-least-squares (ALS) refinement loop that utilizes GPTQ and a Cholesky-weighted truncated SVD. 
Figure~\ref{fig:scientistone_pipeline_parameter_golf} illustrates the complete pipeline of these novel ideas generated by \sys{}.
Internal ablations isolate the ALS loop as the primary driver of the performance gain. 
In contrast, DeepScientist introduced no algorithmic changes.
Its modifications were limited to environment and portability adjustments. 
As a result, DeepScientist failed to improve the underlying algorithm -- merely replicating the reference's performance -- and ultimately produced an invalid submission by exceeding the strict 16MB artifact size limit. 
% Full details are available in our generated paper provided in the supplementary materials.
%In short, only \sys{} delivered a valid, state-of-the-art solution driven by a substantive, structurally novel algorithmic delta. 

% \paragraph{Idea-evolution}

%\todo{Stats: Adapt paper\_stats.py and add results here.}


% We highlight two medal-achieving runs for MLEBench tasks, each illustrating a \emph{Parallel Explore-Exploit} discovery trajectory.


% \paragraph{Case 1: Conservative exploitation (iMet-2020, Silver).}
% The i1 explore pool produced four diverse hypotheses with a best score of $0.588$ (``Sub-Center Disjoint Box Geometries for Polymorphic Art Lineages'').
% At i2 the conservative critic identified asymmetric-loss formulations as the strongest signal in the elite pool and proposed ``Data-Driven Hierarchical Asymmetric Loss for Noisy Multi-Label Classification'' (i2.b1, $0.633$); the reflective checkpointer marked it for exploit.
% Conservative refinements of the same idea reached $0.674$ at i3.b1 and $0.679$ at i4.b1---the silver-medal score---a textbook explore-then-exploit trajectory in which three rounds of incremental refinement compounded to a $7$-point gain over the best i1 explore.

% \paragraph{Case 2: Late-stage breakthrough from a modest signal (3D object detection, Gold).}
% No branch exceeded score of $0.122$ through the first three iterations (against a $0.042$ median---already bronze-clearing, but far from gold).
% The i3.b2 explore branch ``Unified Semantic Positional Encoding'' scored only $0.085$ but was flagged by the checkpointer as a promising under-explored axis.
% At i4 the conservative critic picked up that signal and proposed ``Hybrid Semantic-Cartesian Positional Encodings'' (i4.b2), which scored $0.176$---more than doubling the previous best and clearing the gold threshold.
% The journal makes the chain visible: a sub-median explore branch became the directional signal that the next iteration's conservative critic refined into a gold-medal result.




%\paragraph{\coeaudit{} results}
%Do the same failure modes (hallucinated references, method-code misalignment, score inflation) appear on non-ADRS tasks? Does CoE Audit generalize?


% \todo{Present results beyond ADRS to test whether \coeaudit{} and the failure modes observed in \S\ref{sec:experiments} generalize to other benchmarks.

% Comparison: \sys{} vs DeepScientist on selected MLE-Bench and ParameterGolf tasks. Show:
% \begin{enumerate}
% \item Solver performance (if applicable).
% \item \coeaudit{} Link Check results (same four checks applied identically).
% \item Key finding: do the same failure modes (hallucinated references, method-code misalignment, score inflation) appear on non-ADRS tasks? Does \coeaudit{} generalize?
% \end{enumerate}

% This section should be a compact table + 2-3 paragraphs.
% The main claim is that CoE and \coeaudit{} are benchmark-agnostic: the evidence chain framework and Link Checks apply wherever deterministic evaluators and code artifacts exist.}

